\chapter{Теоретические сведения}
\label{ch:theory}

\section{JavaScript как язык программирования}

\textbf{JavaScript} (часто сокращается до \textbf{JS}) --- интерпретируемый
язык программирования с динамической типизацией, разработанный в
\textit{Netscape Communications} в 1995~году Бренданом Айком. Изначально
создавался для добавления интерактивности в веб-страницы, в дальнейшем
стандартизирован как \textbf{ECMAScript} (текущая редакция --- ES2024).

Ключевые особенности языка:
\begin{itemize}
  \item \textbf{Интерпретируемость} --- код выполняется построчно
        средой исполнения (браузер, \texttt{Node.js}), без отдельного
        этапа компиляции в машинный код;
  \item \textbf{Динамическая типизация} --- тип переменной определяется
        значением во время выполнения и может меняться;
  \item \textbf{Прототипная объектная модель} --- наследование строится
        через цепочки прототипов, а не через классы (синтаксис
        \texttt{class} с ES2015 --- лишь синтаксический сахар над
        прототипами);
  \item \textbf{Функции первого класса} --- функции являются полноценными
        значениями: их можно сохранять в переменных, передавать в
        аргументы, возвращать из других функций;
  \item \textbf{Однопоточная асинхронная модель} --- ровно один поток
        исполнения и очередь задач (\textit{event loop}); долгие
        операции (сетевые запросы, таймеры) выполняются вне основного
        потока и возвращают результат через колбэки или промисы.
\end{itemize}

\section{Переменные и типы данных}

Объявление переменных в современном JavaScript --- через ключевые слова
\texttt{let} (изменяемая) и \texttt{const} (неизменяемая ссылка).
Устаревшее \texttt{var} не используется: оно не поддерживает блочную
область видимости и приводит к скрытым ошибкам.

\begin{lstlisting}[caption={Объявление переменных и базовые типы},captionpos=b]
const username = 'my name';   // строка, неизменяемая
let bonusBalance = 1000;      // число, изменяемая
const isActive = true;        // логический тип
const messages = [];          // массив (объект)
const user = { name: 'Ann' }; // объект
let absent;                   // undefined
const empty = null;           // null
\end{lstlisting}

Семь примитивных типов: \texttt{string}, \texttt{number},
\texttt{boolean}, \texttt{null}, \texttt{undefined}, \texttt{symbol},
\texttt{bigint}. Всё остальное (массивы, функции, даты, регулярные
выражения, экземпляры классов) --- объекты.

\section{Массивы и их методы}

Массив в JavaScript --- индексируемая коллекция произвольных значений.
Базовые методы:

\begin{itemize}
  \item \texttt{push(x)} / \texttt{pop()} --- добавить/удалить элемент с
        конца;
  \item \texttt{forEach(fn)} --- вызвать функцию для каждого элемента;
  \item \texttt{map(fn)} --- получить новый массив, применив функцию к
        каждому элементу;
  \item \texttt{filter(fn)} --- получить новый массив только из тех
        элементов, для которых функция вернула \texttt{true};
  \item \texttt{join(sep)} --- объединить элементы в строку через
        разделитель;
  \item \texttt{includes(value)} --- проверить, есть ли значение в
        массиве (\texttt{true}/\texttt{false}).
\end{itemize}

У строк есть одноимённый метод \texttt{includes}, который проверяет
наличие подстроки:

\begin{lstlisting}[caption={Поиск подстроки методом \texttt{includes}},captionpos=b]
'Чёрный чай'.includes('зелёный'); // false
'Кажется, пойдём в кино'.includes('кино'); // true
\end{lstlisting}

\section{Управляющие конструкции}

Условный оператор \texttt{if/else~if/else}, тернарный оператор
\texttt{a~?~b~:~c}, оператор выбора \texttt{switch}. Циклы:
\texttt{for} (классический счётный), \texttt{for...of} (по значениям
итерируемого объекта), \texttt{for...in} (по ключам объекта),
\texttt{while} и \texttt{do~while}.

\begin{lstlisting}[caption={Цикл с условным выбором имени отправителя},captionpos=b]
for (let i = 0; i < messages.length; i++) {
    const author = i % 2 === 0 ? 'Друг' : 'Вы';
    console.log(author + ': ' + messages[i]);
}
\end{lstlisting}

\section{Функции}

Функция в JavaScript --- объект, который можно вызывать. Способы
объявления:

\begin{lstlisting}[caption={Три способа объявить функцию},captionpos=b]
// 1. Декларация функции
function sum(a, b) { return a + b; }

// 2. Функциональное выражение
const sum = function(a, b) { return a + b; };

// 3. Стрелочная функция (ES2015)
const sum = (a, b) => a + b;
\end{lstlisting}

Стрелочные функции компактны и не имеют собственного \texttt{this} ---
удобно для колбэков. Декларации <<всплывают>> (\textit{hoisting}) ---
их можно вызывать выше по тексту, чем сама декларация. Функциональные
выражения и стрелочные функции \textbf{не} всплывают.

\section{Объекты и работа с JSON}

Объект --- неупорядоченное множество пар <<ключ--значение>>. Запись
объекта в текстовом виде --- \textbf{JSON} (\textit{JavaScript Object
Notation}) --- де-факто стандарт обмена данными в вебе. Преобразование
между JSON-строкой и объектом:

\begin{lstlisting}[caption={\texttt{JSON.parse} и \texttt{JSON.stringify}},captionpos=b]
const obj = JSON.parse('{"name": "Spider-Man", "id": 1009610}');
const str = JSON.stringify(obj); // '{"name":"Spider-Man","id":1009610}'
\end{lstlisting}

\section{DOM API}

\textbf{DOM} (\textit{Document Object Model}) --- объектная модель
HTML-документа: страница представлена деревом узлов
(\texttt{Element}, \texttt{Text}, \texttt{Comment}). Через JavaScript
доступны методы поиска узлов и их изменения:

\begin{itemize}
  \item \texttt{document.getElementById(id)} --- узел по атрибуту
        \texttt{id};
  \item \texttt{document.querySelector(selector)} --- первый узел по
        CSS-селектору; \texttt{querySelectorAll} --- все совпадения;
  \item \texttt{element.innerHTML = '...'} --- заменить внутреннее
        содержимое узла на готовую HTML-строку (быстро, но требует
        самостоятельного экранирования пользовательских данных);
  \item \texttt{element.textContent = '...'} --- задать текст без
        интерпретации тегов (безопасно по умолчанию);
  \item \texttt{element.addEventListener(event, handler)} ---
        подписаться на событие (\texttt{click}, \texttt{submit},
        \texttt{DOMContentLoaded} и др.).
\end{itemize}

Событие \texttt{DOMContentLoaded} срабатывает, когда браузер построил
DOM-дерево (но ещё не загрузил картинки и стили). Это правильный
момент для запуска скриптов, которые ищут элементы по \texttt{id}.

\section{Асинхронность: \texttt{Promise} и \texttt{async/await}}

Сетевые запросы выполняются в фоне и завершаются непредсказуемо ---
JavaScript описывает такой результат объектом \textbf{Promise}.
Промис находится в одном из трёх состояний: \textit{pending} (ожидание),
\textit{fulfilled} (успех) или \textit{rejected} (ошибка). Цепочки
\texttt{.then(...)}\,/\,\texttt{.catch(...)} в современном коде заменены
на синтаксис \texttt{async/await}, который делает асинхронный код
визуально похожим на синхронный.

Запрос к серверу выполняется через встроенный \texttt{fetch}:

\begin{lstlisting}[caption={Асинхронный запрос с обработкой ошибки},captionpos=b]
async function fetchCharacters() {
    const response = await fetch('https://jsfree-les-3-api.onrender.com/characters');
    if (!response.ok) {
        throw new Error('HTTP ' + response.status);
    }
    return await response.json();
}
\end{lstlisting}

\texttt{fetch} возвращает промис, который разрешается объектом
\texttt{Response}; чтобы получить распарсенный JSON, нужен ещё один
\texttt{await response.json()}. Ошибки сети попадают в
\texttt{catch}, ошибки уровня HTTP (\(4xx\), \(5xx\)) --- нет:
\texttt{response.ok} проверяется вручную.

\section{Подключение скриптов в HTML}

JavaScript-файлы подключаются тегом \texttt{<script>} в \texttt{<head>}
или в конце \texttt{<body>}. По умолчанию браузер блокирует разбор
HTML на время загрузки и выполнения скрипта --- поэтому скрипты,
которые работают с DOM, удобно размещать в конце \texttt{<body>}.
Атрибуты \texttt{defer} и \texttt{async} позволяют загружать скрипты
параллельно с разбором HTML; \texttt{defer} дополнительно гарантирует
порядок выполнения и запуск после построения DOM. В учебном проекте
выбран простой вариант --- три тега \texttt{<script>} в конце
\texttt{<body>} (Bootstrap-бандл, \texttt{index.js}, \texttt{start.js}).

\endinput
